\section{Preliminaries}\label{sec:prelim}

\subsection{Notation}

Throughout, $n$ is the field size parameter, $q=2^n$, and $\Fq$ is the finite field with $q$
elements, realised as $\F_2[X]/(f)$ for the irreducible polynomial $f$ fixed by the AIMer
specification (Fact~1.4 of \texttt{docs/FACT\_SHEET.md}: $X^{128}+X^7+X^2+X+1$,
$X^{192}+X^7+X^2+X+1$, $X^{256}+X^{10}+X^5+X^2+1$). A field element $a=\sum_{i<n} a_i X^i$ is
identified with the bit vector $(a_0,\dots,a_{n-1})$; this is the specification's own convention
(\S4.1.1, Fact~0.1). We write $\Tr:\Fq\to\F_2$ for the absolute trace, $\wt(k)$ for the Hamming
weight of the binary expansion of an integer $k$, and $\varphi:a\mapsto a^2$ for the Frobenius
automorphism.

\subsection{AIM3}

\begin{definition}[AIM3, {\cite[\S3.1]{AIMerV3}}; Facts~1.1--1.3 of the fact sheet]\label{def:aim3}
Fix $\ell\ge 1$ and exponents $e_0,\dots,e_\ell$ with $\gcd(e_i,n)=1$. The \emph{enhanced Mersenne
S-box} is
\[
  \Sbox{e}(x) \;=\; x^{2^e-1}+x^{2^n-2},\qquad \Sbox{e}(0)=0 ,
\]
and, writing $\mathrm{Lin}_{in}[iv](x)=(A_0x+b_0)\|\cdots\|(A_{\ell-1}x+b_{\ell-1})$ and
$\mathrm{Lin}[iv](x_0,\dots,x_{\ell-1})=b_\ell+\sum_i A_{\ell+i}x_i$ for $\F_2$-linear maps $A_j$
and constants $b_j$ derived from $iv$,
\[
  \mathrm{AIM3}(pt,iv)\;=\;\Sbox{e_\ell}\circ \mathrm{Lin}[iv]\circ
    \Sbox{e_0,\dots,e_{\ell-1}}\circ \mathrm{Lin}_{in}[iv](pt)\;\oplus\;pt .
\]
Key generation rejects any $(pt,iv)$ for which some S-box input is zero (Algorithm~5, Fact~1.7), so
every S-box relation below may be used on nonzero inputs only.
\end{definition}

The $A_j$ are produced from a $\mathrm{SHAKE}_\lambda(iv)$ tape as $A_j=L_jU_j$ with $L_j$ unit lower
and $U_j$ unit upper triangular (Algorithm~4, Fact~1.6). Two encoding details that the specification
text does not fix unambiguously --- the bit order inside a tape byte, and whether the stored matrix
is $A$ or $A^{\mathsf T}$ --- were pinned against the official known-answer tests; a third axis (the
tape matrix is filled column-wise, and $A=L\cdot U$ rather than $U\cdot L$) had to be added before a
match was possible. The resulting convention reproduces the KAT for all three parameter sets.
%% ledger: O
\begin{remark}\label{rem:conventions}
All \emph{toy}-$n$ experiments in this paper are convention independent: they use random unit-LU
matrices from a seeded PRNG, never the $iv$-derived ones. The convention matters only for statements
about the recommended parameter sets.
\end{remark}

We write $t_i=\Sbox{e_i}(x_i)$ for the S-box outputs and $t_\ell=pt\oplus ct$. The system the
designers analyse is
\[
  \Squad:\quad \text{for each S-box } i:\quad
  P_1^{(i)}=x_iy_i+x_i^{2^{e_i}}+1,\quad P_2^{(i)}=x_iP_1^{(i)},\quad P_3^{(i)}=y_iP_1^{(i)} ,
\]
with $(y_i)=(t_i)$, giving $3(\ell+1)n$ Boolean quadratic equations in $(\ell+1)n$ variables.

\subsection{Boolean ideals, Macaulay matrices and degree falls}

Let $\Bool_N=\F_2[z_1,\dots,z_N]/(z_i^2-z_i)$ be the Boolean polynomial ring; its monomials are the
squarefree ones, and the \emph{Boolean degree} of a monomial is the size of its support. A field
polynomial $F\in\Fq[x,y]$ is turned into $n$ Boolean polynomials by taking coordinates; equivalently,
the $\F_2$-span of its coordinates is $\{\Tr(\gamma F):\gamma\in\Fq\}$, an $n$-dimensional space. We
use the second description throughout, because it makes Frobenius cancellations visible:
\begin{equation}\label{eq:trace-frobenius}
  \Tr\bigl(\gamma\, z^{2^s}\bigr)=\Tr\bigl(\gamma^{2^{-s}} z\bigr),
\end{equation}
so two trace terms coincide exactly when their exponents lie in one Frobenius orbit.

\begin{definition}[Macaulay matrix, degree fall]\label{def:macaulay}
For a system $\mathcal E\subset\Bool_N$ of quadratics and $D\ge2$, the degree-$D$ Macaulay matrix
$\Mac_D(\mathcal E)$ has one row for each product $\mu\cdot f$ with $f\in\mathcal E$ and $\mu$ a
squarefree monomial of degree at most $D-2$, expanded in the monomial basis of degree $\le D$. We
write $\syz_D=\#\text{rows}-\operatorname{rank}\Mac_D$ and
\[
  \mathrm{low}_D=\dim\bigl(\operatorname{rowspace}\Mac_D\cap \Bool_N^{\le D-1}\bigr),\qquad
  \mathrm{falls}_D=\mathrm{low}_D-\operatorname{rank}\Mac_{D-1} .
\]
The \emph{first-fall degree} is the smallest $D$ with $\mathrm{falls}_D>0$.
\end{definition}

\begin{remark}[saturation]\label{rem:saturation}
$\mathrm{falls}_D$ is only meaningful while $\mathrm{low}_D$ is far from $\binom{N}{\le D-1}$. When a
system is strongly over-determined the row space fills the whole low-degree space and
$\mathrm{falls}_D$ measures the monomial basis rather than the algebra. Every fall count reported
here is accompanied by its saturation ratio, and a control system of the same shape is measured
alongside; see Section~\ref{sec:ideal}.
\end{remark}

\subsection{The two semi-regular baselines}

The designers estimate the cost of solving $\Squad$ from a semi-regular Hilbert series. Two different
series are in circulation, and they do not agree.
%% ledger: J
\begin{proposition}[reproduction of the official table]\label{prop:baselines}
\Proved\ (arithmetic) With $N=(\ell+1)n$ variables and $m=3(\ell+1)n$ quadratics, the
\emph{general-field} series $(1-z^2)^m/(1-z)^{N-k}$ together with the cost
$2^k\binom{N-k+d}{d}^2$ reproduces the published figures $262.6/381.0/651.5$ exactly (to $0.1$ bit)
at $k=57/102/150$, where $d_{\mathrm{reg}}=19/25/44$. The \emph{Boolean} series
$(1+z)^N/\bigl((1-z)(1+z^2)^m\bigr)$, which is the one written in the specification, gives
$218.4/312.5/529.1$ instead.
\end{proposition}

We therefore fix two baselines, $B_1$ (the published numbers) and $B_2$ (the Boolean series), and a
third, $B_3$, the cost of exhaustive search. The conversion between AIM3 evaluations and bit
operations is fixed in \texttt{docs/COMPLEXITY\_ACCOUNTING.md} from the addition chains of the
specification's Table~5. No attack complexity is claimed anywhere in this paper;
Proposition~\ref{prop:baselines} is a calibration statement, not a correction of the designers.
